%% --NOTAÇÃO PADRÃO--

%> Alfabetos matemáticos e operadores
%\DeclareMathAlphabet{\pzcal}{OT1}{pzc}{m}{it} %> Usar \pzcal
%\DeclareMathAlphabet{\dcal}{U}{dutchcal}{m}{n} %> Usar \dcal
%\DeclareMathAlphabet{\pzcal}{OMS}{zplm}{m}{n}
%\DeclareMathOperator{\Tr}{Tr}

\setmathfont[range={scr, bfscr}]{XITS Math}
\DeclareMathAlphabet{\pzcal}{OMS}{zplm}{m}{n}
\let\dcal\mathscr
\newcommand{\upmu}{\symup{\mu}}
\DeclareMathOperator{\Tr}{Tr}


%> Notação para variáveis e grandezas
%> Diferencial
\newcommand{\diff}{
    \mathop{}\!\symup{d}}

%> Grandezas caligráficas por partícula
% NENHUMA MUDANÇA AQUI! Como preservamos o \dcal no passo anterior, 
% isso vai continuar funcionando perfeitamente.
\newcommand{\pp}[1]{ %per particle [X]
    \dcal{\MakeLowercase{#1}}}
    
\newcommand{\pps}{ %per particle entropy
    \dcal{s}}
    
\newcommand{\ppu}{ %per particle energy
    \dcal{u}}
  
\newcommand{\zc}{
    \pzcal{Z}_C
}

\newcommand{\zg}{
    \pzcal{Z}_G
}

%> Coordenadas e momentos (Slanted text)
% Substituímos a chamada de baixo nível (usefont) pela chamada de alto nível do LaTeX.
% \normalfont garante a fonte principal do texto, e \slshape garante a inclinação "slanted".
\newcommand{\p}{
  \text{\normalfont\slshape p}}
  
\newcommand{\q}{
  \text{\normalfont\slshape q}}
  
\newcommand{\R}{
  \text{\normalfont\slshape r}}
  
\newcommand{\X}{
  \text{\normalfont\slshape x}}
